\ssr{ВВЕДЕНИЕ}

К числу проблем оценки эффективности труда персонала относятся формальность оценки профессиональной деятельности, сложность оценки в случаях, когда результаты являются результатами совместной с кем-то деятельности, либо, когда между каким-то видом деятельности и его результатами проходит достаточно большой промежуток времени. Также к трудностям такой оценки можно отнести недопонимание получаемых коэффициентов и показателей эффективности деятельности персонала и их влияния на стимулирование труда.

Существующие методы оценки результативности трудовой деятельности персонала в настоящее время далеки от совершенства. Их влияние на показатели результативности имеет двоякий и не всегда положительный характер, а мотивация сотрудников не всегда учитывается~\cite{Ilchenko}.

Цель работы~---~разработка базы данных для хранения и обработки оценок трудовой эффективности сотрудников группы компаний.

Для достижения поставленной цели необходимо решить следующие задачи:
\begin{enumerate}[label={\arabic*)}]
	\item провести анализ предметной области, выделить основные сущности и связи между ними, сформулировать требования и ограничения к разрабатываемой базе данных;
	\item спроектировать сущности базы данных и определить ограничения целостности данных, спроектировать ролевую модель на уровне базы данных;
	\item выбрать средства реализации базы данных и приложения;
	\item реализовать сущности базы данных и ограничения целостности. Описать методы тестирования и разработать тестовые случаи для проверки корректности работы функционала;
	\item провести исследование производительности базы данных при увеличении объема данных и количестве одновременно выполняемых запросов, а также с использованием кеширования и без него.
\end{enumerate}
\clearpage	